\section{A proved local part of the normalized-radius conjecture}
\label{sec:local-radius}

The earlier \emph{Angular Zero Continuation} report \cite{AngularBaseline}
asks whether
\[
 \eta_{a,b}(\rho)=\frac{\cos\theta_{a,b}(\rho)}{\rho}
\]
decreases on the whole interval $0<\rho\le1$ for every integer
$a\ge2$, $b>0$. At $a=1$ it predicts decreasing behavior for $b>1$,
increasing behavior for $0<b<1$, and the constant $1/2$ for $b=1$.
That report identifies the sign of the first nonconstant Taylor
coefficient as a smaller explicit open problem. We solve that problem
for the full stated parameter range and describe the intervening
fractional-order threshold.

Throughout this section put
\[
 U=1+2^{-b},\qquad V=1+2^{-b}+3^{-b},\qquad
 W=1+2^{-b}+3^{-b}+4^{-b}.
\]
The parameter $\rho$ in this section is a disk radius. It is unrelated
to the Lerch deformation parameter used later.

\begin{lemma}[The local radial coefficient]\label{lem:radius-coefficient}
For fixed $a\ge1$, $b>0$, $\eta_{a,b}$ extends as a real analytic
even function near $\rho=0$, and
\begin{equation}\label{eq:eta-expansion}
 \eta_{a,b}(\rho)=\frac U2\left(\frac23\right)^a
                    +K(a,b)\rho^2+O_{a,b}(\rho^4),
\end{equation}
where
\begin{equation}\label{eq:eta-K}
 K(a,b)=UV\,3^{-a}-\frac W2\left(\frac25\right)^a
                  -\frac{U^3}{2}\left(\frac8{27}\right)^a.
\end{equation}
\end{lemma}

\begin{proof}
Set $\delta=\pi/2-\theta$ and divide
$\operatorname{Im}F_{a,b}(\rho e^{i\theta})$ by $2^{-a}\rho^2$.
The resulting equation is analytic near $(\rho,\delta)=(0,0)$
and starts with $\sin(2\delta)$; its $\delta$ derivative there is
$2$. The implicit-function theorem gives a unique analytic solution,
odd under $(\rho,\delta)\mapsto(-\rho,-\delta)$.
Writing $\delta=\alpha\rho+\beta\rho^3+O(\rho^5)$, the modes
with outer indices $2,3,4,5$ give
\[
 \alpha=\frac U2(2/3)^a,\qquad
 \beta=UV\,3^{-a}-\frac W2(2/5)^a
                   -\frac{23U^3}{48}(8/27)^a.
\]
All higher modes enter only at order $\rho^5$ or later in the
normalized equation. Since $\eta=\sin\delta/\rho$, its quadratic
coefficient is $\beta-\alpha^3/6$, which is \eqref{eq:eta-K}.
This also rederives the earlier report's local coefficient without
assuming its sign.
\end{proof}

\begin{theorem}[Complete sign in the conjectured local range]
\label{thm:radius-local-sign}
The coefficient \eqref{eq:eta-K} satisfies
\begin{enumerate}
\item $K(a,b)<0$ for every real $a\ge2$ and every $b>0$;
\item $K(1,b)<0$ if $b>1$, $K(1,b)>0$ if $0<b<1$, and
      $K(1,1)=0$;
\item for $b\ge1$, $K(a,b)<0$ whenever $a>1$;
\item for each $0<b<1$, there is exactly one $a_*(b)\in(1,2)$
      with $K(a_*(b),b)=0$. The coefficient is positive on
      $1\le a<a_*(b)$ and negative for $a>a_*(b)$.
\end{enumerate}
The function $a_*(b)$ is real analytic on $(0,1)$. Every strict
coefficient sign gives the corresponding strict monotonicity of
$\eta_{a,b}$ on some interval $0<\rho<\rho_0(a,b)$. No all-radius
monotonicity claim is made here.
\end{theorem}

\subsection{Monotonicity in the outer order}
Define
\[
 G_a=W+U^3\left(\frac{20}{27}\right)^a
            -2UV\left(\frac56\right)^a.
\]
Then $K(a,b)=-\tfrac12(2/5)^aG_a$, so it suffices to determine
the sign of $G_a$. We first show
\begin{equation}\label{eq:G-increasing}
 \frac{\partial G_a}{\partial a}>0\qquad(a\ge1, b>0).
\end{equation}
Put $x=2^{-b}$ and $y=3^{-b}$. Since $y\ge x^2$ and $0<x<1$,
\[
 \frac{U^2}{V}\le\frac{(1+x)^2}{1+x+x^2}\le\frac43.
\]
The last inequality is equivalent to $(1-x)^2\ge0$.
Elementary alternating-series bounds for $\log(1+t)$ give
\[
 \log(27/20)<\frac{7273}{24000},\qquad
 \log(6/5)>\frac9{50}.
\]
Consequently
\[
 \frac{U^2}{2V}
       \frac{\log(27/20)}{\log(6/5)}
 <\frac{7273}{6480}<\frac98
 \le\left(\frac98\right)^a.
\]
This is precisely the inequality obtained by differentiating $G_a$
and moving its negative term to the other side. It proves
\eqref{eq:G-increasing}.

\subsection{A finite positive certificate at \texorpdfstring{$a=2$}{a = 2}}
Because $\log_2 3>3/2$, $y\le x^{3/2}$. The coefficient of $y$
in $G_2$ is $-(7+25x)/18<0$. With $x=t^2$, $0<t<1$, we obtain
\begin{equation}\label{eq:G2-polynomial}
 G_2\ge P_6(t):=
 \frac{800t^6-2025t^5+1833t^4-567t^3-192t^2+233}{1458}.
\end{equation}
This polynomial has the degree-six Bernstein representation
\[
 P_6(t)=\sum_{j=0}^6 b_j\binom6j t^j(1-t)^{6-j},
\]
with the exact coefficient vector
\begin{equation}\label{eq:bernstein-six}
 (b_0,\ldots,b_6)=
 \left(\frac{233}{1458},\frac{233}{1458},\frac{367}{2430},
 \frac{665}{5832},\frac{55}{486},\frac{95}{1458},\frac{41}{729}\right).
\end{equation}
Every coefficient is positive, and the Bernstein basis is nonnegative
and sums to one on $[0,1]$. Hence $G_2>0$. Together with
\eqref{eq:G-increasing}, this proves the first assertion of the theorem.
The accompanying verifier expands \eqref{eq:bernstein-six} over
$\mathbb Q$ and checks equality with \eqref{eq:G2-polynomial}; no
floating-point polynomial sign test is used.

\subsection{The sign change at \texorpdfstring{$a=1$, $b=1$}{a = 1, b = 1}}
A direct simplification gives
\begin{equation}\label{eq:G1}
 27G_1=20x^3+42x^2-3x+2-(45x+18)y.
\end{equation}
If $b>1$, let $t=\sqrt{2x}\in(0,1)$. Since
$y=(2x)^{\log_2 3}/3\le t^3/3$, substitution in
\eqref{eq:G1} yields
\[
 G_1\ge(1-t)P_5(t),\qquad
 P_5(t)=\frac{-5t^5+10t^4-11t^3+t^2+4t+4}{54}.
\]
The degree-five Bernstein coefficients of $P_5$ are
\begin{equation}\label{eq:bernstein-five}
 \left(\frac2{27},\frac4{45},\frac{19}{180},
       \frac{14}{135},\frac1{10},\frac1{18}\right).
\end{equation}
They are all positive, so $G_1>0$ for $b>1$.

Suppose instead $0<b<1$, so $1/2<x<1$. Put $p=\log_2 3$.
The exact integer inequality $2^{11}<3^7$ gives $p>11/7$, while
$p<2$. For $y(x)=x^p$ this implies
\[
 y(1/2)=\frac13,\quad y'(1/2)>\frac{22}{21},\quad
 y''(x)=p(p-1)x^{p-2}>\frac{44}{49}\quad(1/2<x<1).
\]
Taylor's formula with an integral remainder therefore gives
\[
 y\ge\frac13+\frac{22}{21}(x-1/2)
                   +\frac{22}{49}(x-1/2)^2.
\]
Using the negative coefficient of $y$ in \eqref{eq:G1}, we find
\[
 G_1\le-\frac{(2x-1)(10x^2-337x+334)}{2646}<0.
\]
Indeed the quadratic is decreasing on $[1/2,1]$ and its value at
one is $7>0$. At $b=1$, direct substitution $x=1/2$, $y=1/3$
gives $G_1=0$. The exact identity
$F_{1,1}(z)=\tfrac12\log^2(1-z)$ gives
$\eta_{1,1}(\rho)=1/2$ for every $0<\rho\le1$.

Finally, $G_a$ is strictly increasing in $a\ge1$. The signs of
$G_1$ and $G_2$ give all the remaining assertions and the unique
threshold. Its analyticity follows from the analytic implicit-function
theorem and \eqref{eq:G-increasing}. Differentiating
\eqref{eq:eta-expansion} shows that a nonzero $K$ controls the sign
of $\eta'$ for sufficiently small positive $\rho$.

\begin{remark}[What remains of the conjecture]
This proves the local sign statement explicitly requested in the earlier
report, including all real $a\ge2$. It does not show that a derivative
cannot change sign farther from the origin. At the fractional threshold
$a=a_*(b)$, the quadratic coefficient vanishes; its quartic successor is
needed even for local monotonicity. The next section proves opposite
quartic signs at two inner orders, producing both maxima and minima.
The full-radius integer conjecture and the general quartic-sign
classification remain distinct unresolved questions.
\end{remark}
